% =============================================================================
%  3.2 Clasificación Garibelt
% =============================================================================

\subsection{Clasificación Garibelt}

% ---
\begin{frame}{Espectro MJG}
  \begingroup
    \centering
    \includegraphics[width=\linewidth,height=0.63\textheight,keepaspectratio]
    {Figures/Tableau/Dashboard_Espectro_Red_Actual.pdf}\par
  \endgroup

  \footnotesize
  \begin{itemize}
    \item Inspirado en \cite{kujala2018}, normaliza cada dimensión en cuatro
          intervalos de igual amplitud: crítico $[0.00, 0.25)$,
          débil $[0.25, 0.50)$, aceptable $[0.50, 0.75)$
          e idóneo $[0.75, 1.00]$.
  \end{itemize}
\end{frame}

% ---
\begingroup
\setbeamertemplate{navigation symbols}{}
\begin{frame}{Dashboard comparativo (Espectro MJG)}
  \centering
  \includegraphics[
    width=\linewidth,
    height=0.80\textheight,
    keepaspectratio,
    trim=23bp 20bp 23bp 18bp,
    clip
  ]{Figures/Tableau/Dashboard_Clasificacion_Red_Actual.pdf}\par
\end{frame}
\endgroup
